\chapter{L'indice di vitalità}
\label{cap:vitalita}

\begin{quote}\small
Questo capitolo risponde all'obiettivo O6, formulato dal relatore come segue:
\emph{oltre alla presenza sì/no, creare un indice di vitalità --- un algoritmo che
valuti quanto una persona si muove: questa persona si muove 50 su 100, quindi è viva
o moderatamente viva}. È l'obiettivo più avanzato della tesi ed è anche la
dimostrazione costruttiva che i dati del radar abilitano ciò che con un bit non è
rappresentabile. Alla data di questa stesura l'algoritmo è specificato e il percorso
di taratura è definito; la taratura sui dati è in corso.
\end{quote}

\section{Motivazione: dalla presenza al triage}
\label{sec:vitalita-motivazione}

Nel contesto del progetto, la differenza fra ``presenza'' e ``vitalità'' è la
differenza fra due domande distinte dei soccorritori:

\begin{enumerate}[leftmargin=1.8em,itemsep=2pt]
  \item \emph{C'è qualcuno sotto il banco?} --- è la domanda di presenza, cui
        rispondono gli obiettivi O1--O3.
  \item \emph{In che condizioni si trova?} --- si muove attivamente, quindi è
        cosciente; compie solo micro-movimenti, quindi è ferito o intrappolato ma
        vivo; oppure non dà segni rilevabili.
\end{enumerate}

Il PIR non può nemmeno porsi la seconda domanda: la sua uscita è un bit, e come
mostrato in \cref{sec:interpretazione-pir} quel bit non misura nemmeno la durata del
movimento. Il radar mmWave sì, perché espone l'\textbf{energia riflessa per gate di
distanza}: una misura continua di \emph{quanto} si muove ciò che riflette, non solo
del fatto che si muova.

L'indice di vitalità è la sintesi di questa informazione in un numero fra 0 e 100,
leggibile da un operatore non tecnico, accompagnato da una classificazione a quattro
livelli.

\section{Dati di ingresso}
\label{sec:vitalita-ingressi}

L'algoritmo lavora sui campioni a 5~Hz prodotti dal logger
(\cref{sec:piattaforma}), usando i canali di \cref{tab:vitalita-ingressi}.

\begin{table*}[htbp]
\centering
\caption{Segnali di ingresso dell'indice di vitalità.}
\label{tab:vitalita-ingressi}
\small
\begin{tabularx}{\linewidth}{@{}llX@{}}
\toprule
\textbf{Segnale} & \textbf{Colonna CSV} & \textbf{Che cosa misura} \\
\midrule
Energia del bersaglio in movimento & \texttt{moving\_energy} &
  intensità del movimento rilevato (0--100) \\
Energia per-gate, movimento & \texttt{menergy\_gate0..8} &
  dove avviene il movimento \\
Energia per-gate, stazionario & \texttt{senergy\_gate0..8} &
  riflessione del corpo fermo; oscilla con il respiro \\
Stato di presenza & \texttt{radar\_presence} &
  condizione di guardia: se nulla, l'indice è forzato a zero \\
\bottomrule
\end{tabularx}
\end{table*}

L'osservazione su cui poggia l'algoritmo è la seguente: una persona \textbf{viva ma
immobile} ha energia in movimento prossima a zero, ma l'energia del suo gate
\emph{oscilla} con il respiro nella banda 0,1--0,5~Hz. Una stanza vuota ha energia
costante, a meno del rumore di fondo. La \textbf{variabilità} dell'energia è dunque
il segnale di vita minimo, quello che sta sotto il movimento.

\paragraph{Vincolo importante ereditato dal \cref{cap:moduli}.} La formulazione
iniziale dell'algoritmo usava l'energia \emph{stazionaria} del gate attivo come
sorgente del termine respiratorio. I risultati di \cref{sec:senergy-zero} e
\cref{sec:respiro} impongono una correzione: sotto 1,5~m i canali stazionari per-gate
sono strutturalmente nulli, e a tutte le distanze provate i canali stazionari
risultano saturi o discordanti mentre i canali \emph{in movimento} forniscono stime
concordi e non sature. L'implementazione deve quindi selezionare la sorgente del
termine respiratorio fra i canali in movimento, o comunque scartare i canali saturi
e quelli identicamente nulli con criterio esplicito. È un esempio di come la fase di
caratterizzazione abbia modificato il progetto dell'algoritmo, e non solo confermato
delle attese.

\section{Algoritmo}
\label{sec:vitalita-algoritmo}

L'algoritmo ha due componenti, aggiornate a ogni campione, e una fusione.

\begin{align}
g^{*} &= \arg\max_i \; e_i && \text{gate attivo} \label{eq:gate-attivo}\\
m_k &= \max\!\left(E_{\text{mov}},\; e^{\text{mov}}_{g^{*}}\right) && \text{movimento istantaneo} \\
M_k &= \alpha_m \, m_k + (1-\alpha_m)\, M_{k-1} && \text{EWMA rapida, } \alpha_m = 0{,}1 \\
d_k &= \left| e_{g^{*}}[k] - e_{g^{*}}[k-1] \right| && \text{variazione tick a tick} \\
R_k &= \alpha_r \, d_k + (1-\alpha_r)\, R_{k-1} && \text{EWMA lenta, } \alpha_r = 0{,}02 \\
V_k &= \mathrm{clamp}\!\left(M_k + k_r R_k,\; 0,\; 100\right) && \text{indice di vitalità}
\label{eq:vitalita}
\end{align}

con la condizione di guardia $V_k = 0$ se \texttt{radar\_presence} $= 0$.

\subsection{Razionale delle scelte}

Ogni scelta ha una motivazione, ed è utile esplicitarla perché l'algoritmo deve
girare identico in Python e su microcontrollore.

\begin{description}[leftmargin=1.4em,style=nextline,itemsep=4pt]
  \item[Media esponenziale e non finestra mobile] La EWMA costa $O(1)$ in memoria
    (due variabili in virgola mobile) e gira senza modifiche sull'ESP32; una finestra
    mobile richiederebbe un buffer circolare e darebbe risultati equivalenti. Poiché
    l'algoritmo deve essere lo stesso in Python e in C++ per essere verificabile
    (passo~5 di \cref{sec:webui-piano}), la formulazione più semplice è quella
    giusta.
  \item[Due costanti di tempo diverse] Il movimento deve reagire in circa 2~s, che a
    5~Hz corrisponde a $\alpha_m = 0{,}1$; il respiro è un segnale lento e va
    integrato su circa 10~s ($\alpha_r = 0{,}02$) per emergere dal rumore. Un'unica
    costante non può servire entrambi gli scopi.
  \item[Massimo fra energia del bersaglio e del gate] Le due grandezze talvolta
    divergono, per effetto dei filtri interni del modulo. Prendere il massimo rende
    l'indice \emph{conservativo verso il falso ``nessun segno''}, che è l'errore più
    grave nel dominio applicativo: dichiarare vuoto un banco sotto cui c'è qualcuno.
  \item[Il coefficiente $k_r$] La variazione $d_k$ è numericamente piccola, di
    qualche unità di energia; $k_r$ la riporta nella scala 10--30 attesa per la classe
    ``vitalità bassa''. È il parametro principale da tarare, con un intervallo
    iniziale plausibile fra 5 e 15.
\end{description}

\section{Classificazione}
\label{sec:vitalita-classi}

\begin{table*}[htbp]
\centering
\caption{Classi dell'indice di vitalità e loro significato operativo. Le soglie
riportate sono iniziali e vanno tarate sui dati (\cref{sec:vitalita-taratura}).}
\label{tab:vitalita-classi}
\small
\begin{tabularx}{\linewidth}{@{}llX@{}}
\toprule
\textbf{Indice} & \textbf{Classe} & \textbf{Significato operativo (triage)} \\
\midrule
0--9 & \texttt{nessun\_segno} & nessuna presenza, o nessuna variazione rilevabile \\
10--39 & \texttt{vitalita\_bassa} & vivo: respiro e micro-movimenti, non si muove \\
40--69 & \texttt{moderato} & movimenti limitati ma attivi \\
70--100 & \texttt{attivo} & movimento pieno, persona reattiva \\
\bottomrule
\end{tabularx}
\end{table*}

\paragraph{Avvertenza da riportare anche nella tesi finale.} Le classi sono un
supporto informativo al triage e \textbf{non costituiscono una diagnosi medica}. In
particolare \texttt{nessun\_segno} significa ``il sensore non rileva variazioni'', non
``persona deceduta'': la persona può essere fuori portata, schermata da materiale
metallico, oppure priva di sensi ma viva. La formulazione dell'etichetta, in
un'interfaccia destinata a soccorritori, è essa stessa una scelta di progetto con
implicazioni: un'etichetta che suggerisca una conclusione clinica sarebbe
inappropriata.

\section{Percorso di sviluppo e taratura}
\label{sec:vitalita-taratura}

La regola metodologica adottata è: \textbf{l'algoritmo si sviluppa su PC, sui file
CSV, dove si può iterare in secondi}. Il porting su ESP32 avviene solo a soglie
validate. Il percorso è il seguente.

\begin{enumerate}[leftmargin=1.8em,itemsep=4pt]
  \item \textbf{Raccolta dei quattro scenari}, cinque trial ciascuno:
        \texttt{vitalita\_0} (stanza vuota), \texttt{vitalita\_respiro} (soggetto
        immobile), \texttt{vitalita\_micro} (micro-movimenti: mano, testa, busto),
        \texttt{vitalita\_attivo} (movimento pieno). Distanza $\geq 1{,}5$~m per
        evitare la saturazione (\cref{sec:saturazione}).
  \item \textbf{Prototipo Python} (\texttt{analisi/vitalita\_proto.py}): implementa
        le equazioni \eqref{eq:gate-attivo}--\eqref{eq:vitalita} e produce la serie
        $V(t)$ per ogni trial.
  \item \textbf{Taratura} di $\alpha_m$, $\alpha_r$, $k_r$ e delle tre soglie in modo
        che i quattro scenari cadano nelle quattro classi attese. Il metodo scelto è
        semplice e difendibile: si osservano le distribuzioni di $V$ per scenario
        (diagrammi a scatola) e si collocano le soglie negli intervalli vuoti fra le
        distribuzioni. Se le distribuzioni si sovrappongono, la sovrapposizione è
        essa stessa un risultato da riportare, e indica che quei due scenari non sono
        separabili con questo indice.
  \item \textbf{Validazione su trial non usati per la taratura}: taratura su
        T01--T03, validazione su T04--T05, con matrice di confusione
        scenario/classe. La metrica finale è la percentuale di campioni classificati
        nella classe attesa, per scenario.
  \item \textbf{Porting} in \texttt{vitality.h} con le stesse costanti, e verifica
        dal vivo nella dashboard (passo~5 di \cref{sec:webui-piano}).
\end{enumerate}

Il passo~4 fornisce il numero da riportare nella tesi: \emph{l'indice classifica
correttamente l'$X\,\%$ dei campioni su scenari mai visti in fase di taratura}. La
separazione fra trial di taratura e trial di validazione è ciò che distingue questo
risultato da una descrizione dei dati su cui l'algoritmo è stato costruito.

\todo{completare con i risultati della taratura: boxplot per scenario, valori
scelti dei parametri, matrice di confusione sui trial di validazione}

\section{Casi limite e comportamenti attesi}
\label{sec:vitalita-limiti}

\begin{table*}[htbp]
\centering
\caption{Casi limite dell'indice di vitalità e comportamento previsto.}
\label{tab:vitalita-casilimite}
\small
\begin{tabularx}{\linewidth}{@{}XXX@{}}
\toprule
\textbf{Situazione} & \textbf{Comportamento atteso} & \textbf{Note} \\
\midrule
La persona esce dal campo & indice a 0 entro il timeout del radar ($\sim$5~s) &
  forzatura sulla condizione di presenza \\
La persona si addormenta & da \texttt{attivo} a \texttt{vitalita\_bassa} in
  30--60~s & le medie esponenziali decadono con le loro costanti \\
Ventilatore o tenda in movimento & possibile falso \texttt{moderato} &
  limite da dichiarare; mitigazione tramite le soglie per-gate del radar \\
Due persone & indice riferito al gate dominante &
  coerente con il limite mono-bersaglio del LD2410B \\
Persona oltre 4--5~m & respiro non più rilevabile, quindi sottostima &
  va documentata la distanza massima utile per la classe
  \texttt{vitalita\_bassa} \\
Persona a meno di 1,5~m & canali stazionari per-gate non disponibili &
  usare i canali in movimento (\cref{sec:senergy-zero}) \\
\bottomrule
\end{tabularx}
\end{table*}

\section{Estensioni possibili}

Tre estensioni sono state individuate ma non sono necessarie per l'obiettivo, e
vanno affrontate solo se il tempo lo consente.

\begin{enumerate}[leftmargin=1.8em,itemsep=2pt]
  \item Sostituire il termine respiratorio a media esponenziale con la
        \textbf{potenza in banda 0,1--0,5~Hz} calcolata su una finestra scorrevole di
        30~s, cioè una FFT ridotta eseguita a bordo. È più robusta agli artefatti ma
        computazionalmente più costosa, e richiede un buffer: sarebbe una versione~2.
  \item \textbf{Isteresi sulle classi}: una classe cambia solo dopo $N$ secondi di
        stabilità, per evitare lo sfarfallio dell'etichetta nell'interfaccia.
  \item \textbf{Indicatore di confidenza} affiancato all'indice, funzione della
        stabilità del gate attivo: un indice calcolato su un gate che cambia
        continuamente è meno affidabile e sarebbe onesto dichiararlo.
\end{enumerate}
